% BEGIN CR28_RADIAL_CONSTITUTIVE_EVALUATION
\section{Continuité constitutive et évaluation de la section gravitationnelle}
\label{cr28:sec:radial-evaluation}

La construction longitudinale détermine \(L\) avant d'utiliser l'identification gravitationnelle. La composition des règles R1--R8 conserve la section de vitesse de \eqref{cr28:eq:same-velocity} et l'action de retour préalablement engendrée. Avec
\[
 L=\frac{5759}{23040}\alpha^{16}L_*,
 \qquad \hbar_{\rm ret}=\eta_{\rm ret}S_*,
 \qquad c_{\rm int}=c_*\widehat c,
\]
le résultat de \eqref{g26:eq:g-rigido-es} s'exprime comme
\begin{equation}
 G_{\rm ret}
 =\frac{c_*^3L_*^2}{S_*}
   \left(\frac{5759}{23040}\right)^2
   \frac{\alpha^{32}\widehat c^{\,3}}{\eta_{\rm ret}}
 =\frac{L^2}
 {\hbar_{\rm ret}(\mu_{\rm vac}\varepsilon_{\rm vac})^{3/2}}.
 \label{cr28:eq:g-constitutive-evaluation}
\end{equation}
Les canaux sont les mêmes que ceux de \eqref{eq:vacuum-angular-channels} :
\[
 q_\pm=\exp\!\left[-\frac{\pi}{180}(A_{\rm deg}\pm C^*_{\rm deg})\right],
 \qquad r_\pm=\frac{1-q_\pm^{90}}{1-q_\pm^{120}},
 \qquad \widehat c=(r_+r_-)^{-1}.
\]
Les bases constitutives satisfont \(\varepsilon=\varepsilon_*r_+^2\), \(\mu=\mu_*r_-^2\) et \(c_*=(\mu_*\varepsilon_*)^{-1/2}\). Leur produit prouve \(c_{\rm int}=(\mu\varepsilon)^{-1/2}\). Si la coordonnée dimensionnelle de \(c_{\rm int}\) est déclarée directement, sa base se retrouve comme \(c_*=c_{\rm int}/\widehat c\). Le coefficient constitutif n'est pas multiplié de nouveau : les deux expressions représentent la même section.

Dans la représentation \(L_*=1\,\mathrm m\), \(S_*=10^{-34}\,\mathrm{J\,s}\) et \(c_{\rm int}=299792458\,\mathrm{m\,s^{-1}}\), l'évaluation des lecteurs et de la composition radiale donne
\begin{equation}
 \boxed{
 G_{\rm ret}
 =6.674299659414022706367776189\ldots\times10^{-11}
 \,\mathrm{m^3\,kg^{-1}\,s^{-2}}.}
 \label{cr28:eq:g-evaluation}
\end{equation}
La valeur de comparaison \(G_{\rm CODATA\,2022}=6.67430(15)\times10^{-11}\,
\mathrm{m^3\,kg^{-1}\,s^{-2}}\) conserve son incertitude expérimentale. Dans les mêmes unités, l'écart de l'évaluation à la valeur centrale est \(-3.40585977\ldots\times10^{-18}\), soit environ \(-0.00227057\) incertitudes-types. Les chiffres de l'évaluation correspondent aux lecteurs et aux bases déclarés ; l'incertitude appartient au résultat mesuré. Aucun chiffre de la référence ne sélectionne \(N\), \(r_N\), la règle longitudinale ou la norme radiale.

\subsection{Réalisation chronologique du même rayon construit}
\label{cr28:sec:radial-clock} Avec la durée élémentaire
\begin{equation}
 t_0=\frac{\pi L}{54c_{\rm int}},
 \qquad \ell_0=c_{\rm int}t_0=\frac{\pi L}{54},
 \label{cr28:eq:clock-from-length}
\end{equation}
le calendrier de \(108\) positions satisfait
\[
 T_{108}=108t_0=\frac{2\pi L}{c_{\rm int}},\qquad
 \omega_{108}=\frac{c_{\rm int}}{L},\qquad
 R_{108}=\frac{c_{\rm int}}{\omega_{108}}=L.
\]
Par substitution,
\begin{equation}
 \left(\frac{54}{\pi}\right)^2
 \frac{c_{\rm int}^5t_0^2}{\hbar_a}
 =\frac{c_{\rm int}^3L^2}{\hbar_a}.
 \label{cr28:eq:circular-radial-agreement}
\end{equation}
Cette identité place le sélecteur circulaire conservé ci-dessous dans la même normalisation longitudinale. Son théorème d'unicité conserve la prémisse géométrique qu'il énonce ; la construction précédente spécifie en outre la longueur utilisée dans cette réalisation.

Le caractère positif \(Y=I\) donne simultanément \(R_X=\Lambda_X=L\) et \(M_X=\hbar_a/(c_{\rm int}L)\). Pour un caractère général \(Y>0\), les longueurs sont réciproques : \(R_X=LY\), \(\Lambda_X=LY^{-1}\). Leur produit demeure \(L^2I\) sans que leurs valeurs individuelles doivent coïncider. L'identification des quatre longueurs du cas circulaire correspond donc à un secteur précis de la preuve radiale et ne supprime pas la dépendance au caractère.

La proportionnalité \(R_X=(G/c_{\rm int}^{2})M_X\) est commune à tous les caractères du domaine. Sa conservation sous changement de base, raffinement, minimisation avec mémoire, complément de Schur et normalisation dynamique a été démontrée dans \cref{g26:sec:rigidez-radial-es}. Le porteur à cinq composantes et sa réduction à deux hélicités, développés ensuite, décrivent la représentation tensorielle ; ils ne remplacent pas la démonstration du coefficient radial commun.
% END CR28_RADIAL_CONSTITUTIVE_EVALUATION
